\section{Related Work and Conclusion}
\label{sec:related}

\paragraph{Fast weights and recurrent memory.}
Fast-weight memories and linear-attention recurrences treat recent
observations as state updates
\citep{schmidhuber1992fastweights,ba2016fastweights,
katharopoulos2020transformersrnn,schlag2021fastweight}.  Delta-rule families
add or factor specific control freedoms
\citep{yang2024deltanet,yang2025gateddeltanet,hatamizadeh2026gdn2,
li2026eda}.  We cite these systems for architectural lineage, not as systems
evaluated by the controlled CATENA experiments.

\paragraph{Editing and operator prediction.}
Model and cache editing study localized changes to stored behavior
\citep{mitchell2022mend,meng2022rome,meng2023memit,li2026kveraser}, while
hypernetworks, dynamic filters, and low-rank adaptation provide precedents for
descriptor-conditioned or low-rank operators
\citep{ha2016hypernetworks,jia2016dynamicfilters,hu2022lora}.  Our scope is a
controlled recurrent-memory update algebra rather than a claim of replication
or superiority over those systems.

\paragraph{Functional intervention.}
The completed core uses intervention, transplant, rescue, and scalarization
as mechanism-level tests.  This connects methodologically to causal mediation
and interchange-intervention work
\citep{pearl2001direct,vig2020causalmediation,
geiger2021causalabstractions}, while retaining the narrower functional claim
defined by the controlled protocol.

\paragraph{Conclusion.}
Across controlled finite memories, learned operator families, and repeated
structured transactions, the results connect error to behavioral
reachability and the minimal controller class to rank, basis algebra, and
specific update freedoms.  The central design rule is therefore to match
memory architecture to the geometry and algebra of the update family it must
realize.  The next bounded tests are semantic demand inference and
quality-controlled official backends; neither is inferred from the present
controlled evidence.
